Die Masse bleibt gleich, das Volumen $V = \frac{4}{3}\pi r^3$ wächst mit dem Radius. Mit der relativen Radiusänderung $\varepsilon = \Delta r/r = \przO$ ist der neue Radius $r' = (1+\varepsilon)\,r$ und
\begin{align}
 \rho' &= \frac{m}{V'} = \frac{m}{\frac{4}{3}\pi r^3\left(1+\varepsilon\right)^3} = \frac{\rho}{\left(1+\varepsilon\right)^3}
\end{align}
Die relative Änderung der Dichte ist
\begin{align}
 \frac{\Delta\rho}{\rho} &= \frac{\rho'}{\rho}-1 = \drhoF \\
   &= \frac{1}{\left(1+\przO\right)^3}-1 \\
   &= \drho \approx \result{\drhoP{0}{2}}
\end{align}
Die Dichte nimmt also um etwa das Dreifache der relativen Radiusänderung ab, denn für kleine $\varepsilon$ ist $\left(1+\varepsilon\right)^{-3} \approx 1-3\varepsilon$.
